% !TEX root = ../main.tex
\chapter{Abstract}

Quantum computing represents an emerging paradigm with transformative potential for data science, yet a significant gap persists between the theoretical promises of quantum algorithms and their practical applicability on real-world data. This work presents a systematic empirical comparative study of quantum versus classical algorithms across eight use cases covering supervised classification (QSVM with quantum kernel), regression (QK-SVR), combinatorial optimisation (QAOA, VQE, SA-QUBO) and one-class anomaly detection (OCSVM-ZZ). Experiments are conducted on three dataset types per case when applicable: data with native quantum structure, practice data from SKG Tecnología (mobility incidents and travel times in Bogotá) and public-domain benchmarks.

Results show that quantum advantage is \emph{real but conditional}: QSVM achieves $+38.3$~pp over quantum-native data, while classical algorithms outperform in classification and regression on real-world data. In combinatorial optimisation, QAOA achieves a $-30.3\%$ cost reduction in resource assignment against greedy heuristics. Use Case~7 illustrates the \emph{hybrid quantum-classical paradigm} applied to the NP-hard minimum dominating set problem: the QAOA-COBYLA loop recovers the optimum with probability $\approx 73\%$. Use Case~8 (one-class anomalies, \emph{not} fraud) shows OCSVM-ZZ superior only when novelty is a ZZ phase on the same spatial support (AUC $0.864$ vs IsolationForest $0.553$); on moons/blobs and injected far outliers, IsolationForest wins.

Two original contributions are presented. The first is the \emph{quantumness score} $\mathcal{Q}(D)$, a composite metric based on kernel alignment ($\Delta A$), Frobenius distance ($d_F$) and differential spectral entropy ($\Delta H$), achieving a Spearman correlation of $\rho_S = 0.833$ ($p = 0.005$, $n = 9$) with empirical QSVM advantage. The second is the \emph{QCR Framework} (Quantum-Compatible Representation), a four-step protocol that automatically selects the optimal pre-mapping $T^*$: on Telco Churn it raises $\mathcal{Q}(D)$ from $-0.228$ to $+0.167$ and empirical advantage from $0$~pp to $+4.2$~pp via data re-uploading with \texttt{reps=3}. The practical conclusion is that quantum algorithms should be adopted as \emph{selective, complementary tools} within the enterprise technology stack — not as replacements for classical methods — being particularly promising for combinatorial optimisation with QUBO formulation. All experiments are conducted with Qiskit~2.3.0 on exact statevector simulation.

\vspace{0.5em}
\noindent\textbf{Keywords:} quantum computing, quantum machine learning, QSVM, quantum kernel, QAOA, VQE, combinatorial optimisation, anomaly detection, one-class SVM, variational algorithms, NISQ era, quantum advantage, quantumness score, QCR Framework.

\textbf{Keywords:} quantum computing, quantum advantage, hybrid quantum-classical algorithm, quantumness score, quantum kernel, QAOA, QSVM, OCSVM-ZZ, QCR Framework, minimum dominating set, NISQ
